\chapter{Testing}
\label{chp:testing}

This is the third and last chapter that shows how to move from writing code
that works to writing software. In this particular chapter you will consider how
to write automated tests for your software.

\begin{note}
In this chapter you will cover:
\begin{itemize}
\item 

Assert statements to test code.

\item 

Testing documentation.

\end{itemize}
\end{note}





\section{Tutorial}
\label{sec:testing_tutorial}

In this tutorial you will write code to ensure the correctness of the software
you
have written in Chapter~\ref{chp:modularisation} and~\ref{chp:documentation}
tutorials.

The software for \texttt{absorption.py} is in fact across two separate files:
\begin{itemize}
\item 

\texttt{absorption.py}: the source code. You will check this using \textbf{unit tests}

\item 

\texttt{README.md}: the documentation. You will check this using \textbf{doc tests}.

\end{itemize}


\subsection{Writing tests for code}
\label{\detokenize{building-tools/07-testing/tutorial/main:writing-tests-for-code}}

Recalling the code written in \texttt{absorption.py} in
Section~\ref{sec:modularisation_tutorial},
there are 4 functions that need to be tested:

\begin{itemize}
\item 

\texttt{get\_long\_run\_state}

\item 

\texttt{extract\_Q}

\item 

\texttt{compute\_N}

\item 

\texttt{compute\_t}

\end{itemize}


In the directory that contains \texttt{absorption.py} create a new Python
file called: \texttt{test\_absorption.py}.


Write the following functions to test each of the functions in
\texttt{absorption.py}:

\begin{python}
import numpy as np

import absorption

def test_long_run_state_for_known_number_of_states():
    """
    This tests the `long_run_state` for a small example matrix
    """
    pi = np.array([1, 0, 0])
    P = np.array([[1 / 2, 1 / 4, 1 / 4], [1 / 3, 1 / 3, 1 / 3], [0, 0, 1]])
    pi_after_5_steps = absorption.get_long_run_state(pi=pi, k=5, P=P)
    assert np.array_equal(pi_after_5_steps, pi @ np.linalg.matrix_power(P, 5)), "Did not get expected result for pi after 5 steps"


def test_long_run_state_when_starting_in_absorbing_state():
    """
    This tests the `long_run_state` for a small example matrix.

    In this test we start in the absorbing state, the state vector should not
    change.
    """
    pi = np.array([0, 0, 1])
    P = np.array([[1 / 2, 1 / 4, 1 / 4], [1 / 3, 1 / 3, 1 / 3], [0, 0, 1]])
    pi_after_5_steps = absorption.get_long_run_state(pi=pi, k=5, P=P)
    assert np.array_equal(pi_after_5_steps, pi)


test_long_run_state_for_known_number_of_states()
test_long_run_state_when_starting_in_absorbing_state()
\end{python}


\begin{note}
The two functions you have written do not include a \texttt{return} statement but
instead include an \texttt{assert} statement. An \texttt{assert} is followed by two values
separated by a comma:
\begin{enumerate}

\item 

A boolean that is to be \texttt{True} or \texttt{False}.

\item 

A string that is output if the boolean is `False`.

\end{enumerate}
\end{note}



To run those tests, run the script at the command line:

\begin{cliin}
$ python test_absorption.py
\end{cliin}


When running the tests if everything has been done correctly there will be no
output: the 2 functions have been called and the assertions have
\textbf{passed}. See Figure~\ref{fig:running_tests_with_no_errors}.

\begin{figure}[htbp]
\centering

\includegraphics[width=0.750\linewidth]{./assets/running_tests_with_no_errors/main.png}
\caption{Running the tests with no errors.}
\label{fig:running_tests_with_no_errors}
\end{figure}


For each of the four functions in \texttt{absorption.py} you can now add further tests
and ensure they are also called at the end. The full \texttt{test\_absorption.py} file
should look like:

\begin{python}
import numpy as np

import absorption

def test_long_run_state_for_known_number_of_states():
    """
    This tests the `long_run_state` for a small example matrix
    """
    pi = np.array([1, 0, 0])
    P = np.array([[1 / 2, 1 / 4, 1 / 4], [1 / 3, 1 / 3, 1 / 3], [0, 0, 1]])
    pi_after_5_steps = absorption.get_long_run_state(pi=pi, k=5, P=P)
    assert np.array_equal(pi_after_5_steps, pi @ np.linalg.matrix_power(P, 5)), "Did not get expected result for pi after 5 steps"


def test_long_run_state_when_starting_in_absorbing_state():
    """
    This tests the `long_run_state` for a small example matrix.

    In this test we start in the absorbing state, the state vector should not
    change.
    """
    pi = np.array([0, 0, 1])
    P = np.array([[1 / 2, 1 / 4, 1 / 4], [1 / 3, 1 / 3, 1 / 3], [0, 0, 1]])
    pi_after_5_steps = absorption.get_long_run_state(pi=pi, k=5, P=P)
    assert np.array_equal(pi_after_5_steps, pi)


def test_extract_Q():
    """
    This tests that the submatrix Q can be extracted from a given matrix P.
    """
    P = np.array([[1 / 2, 1 / 4, 1 / 4], [1 / 3, 1 / 3, 1 / 3], [0, 0, 1]])
    Q = absorption.extract_Q(P)
    expected_Q = np.array([[1 / 2, 1 / 4], [1 / 3, 1 / 3]])
    assert np.array_equal(Q, expected_Q), f"The expected Q did not match, the code obtained {Q}"


def test_compute_N():
    """
    This tests the computation of the fundamental matrix N
    """
    P = np.array([[1 / 2, 1 / 4, 1 / 4], [1 / 3, 1 / 3, 1 / 3], [0, 0, 1]])
    Q = absorption.extract_Q(P)
    N = absorption.compute_N(Q)
    expected_N = np.array([[8 / 3, 1], [4 / 3, 2]])
    assert np.allclose(N, expected_N), f"The expected N did not match, the code obtained {N}"


def test_compute_t():
    """
    This tests the computation of the number of steps until absorption t.
    """
    P = np.array([[1 / 2, 1 / 4, 1 / 4], [1 / 3, 1 / 3, 1 / 3], [0, 0, 1]])
    t = absorption.compute_t(P)
    expected_t = np.array([11 / 3, 10 / 3])
    assert np.allclose(t, expected_t), f"The expected t did not match, the code obtained {t}"


test_long_run_state_for_known_number_of_states()
test_long_run_state_when_starting_in_absorbing_state()
test_extract_Q()
test_compute_N()
test_compute_t()
\end{python}


\begin{note}
The \texttt{numpy.array\_equal} and \texttt{numpy.allclose} compare equality of
boolean arrays. They return \texttt{True} or \texttt{False} depending on whether the two
passed arrays are equal or approximately equal (respectively).

\texttt{numpy.allclose} should be used when comparing numpy arrays that might be
different due to numerical imprecision.
\end{note}



\textbf{You can experiment by changing some of the code or the tests and see the way
the tests fail.} See Figure~\ref{fig:running_tests_with_errors_in_source_code}
where the following specific error has been introduced in to \texttt{absorption.py}:
\texttt{P.diagonal() == 1} is incorrect and should be \texttt{P.diagonal() != 1}.

\begin{figure}[htbp]
\centering

\includegraphics[width=0.750\linewidth]{./assets/running_tests_with_error_in_source_code/main.png}
\caption{Running the tests with an error in the source code.}
\label{fig:running_tests_with_errors_in_source_code}
\end{figure}


As and when you add more features to \texttt{absorption.py} you will also add tests.


Software is compromised of both code and documentation. So far you have tested
your code, now you will test your documentation.


\subsection{Testing documentation}
\label{\detokenize{building-tools/07-testing/tutorial/main:testing-documentation}}

To be able to check the python code written in the documentation (see
Chapter~\ref{chp:documentation}) is
correct you need to write the code using a specific format:

\begin{itemize}
\item 

\texttt{>>>} to denote python code

\item 

\texttt{...} to denote secondary lines of multi line python code.

\item 

Nothing to denote the expected output.

\end{itemize}


As an example, in Section~\ref{sec:documentation_tutorial}, you have written:

\begin{md}
```python
import numpy as np

import absorption

pi = np.array([1, 0, 0])
P = np.array([[1 / 2, 1 / 4, 1 / 4], [1 / 3, 1 / 3, 1 / 3], [0, 0, 1]])
```

We now see how the vector $\pi$ changes over time:

```python
for k in range(10):
    print(absorption.get_long_run_state(pi, k, P))
```

This will give:

```
[1. 0. 0.]
[0.5  0.25 0.25]
[0.33333333 0.20833333 0.45833333]
[0.23611111 0.15277778 0.61111111]
[0.16898148 0.1099537  0.72106481]
[0.12114198 0.0788966  0.79996142]
[0.08686986 0.05658436 0.85654578]
[0.06229638 0.04057892 0.8971247 ]
[0.0446745 0.0291004 0.9262251]
[0.03203738 0.02086876 0.94709386]

```
\end{md}


You will modify the above to be:

\begin{md}
```python
>>> import numpy as np
>>> import absorption
>>> pi = np.array([1, 0, 0])
>>> P = np.array([[1 / 2, 1 / 4, 1 / 4], [1 / 3, 1 / 3, 1 / 3], [0, 0, 1]])

```

We now see how the vector $\pi$ changes over time:

```python
>>> for k in range(10):
...     print(absorption.get_long_run_state(pi, k, P))
[1. 0. 0.]
[0.5  0.25 0.25]
[0.33333333 0.20833333 0.45833333]
[0.23611111 0.15277778 0.61111111]
[0.16898148 0.1099537  0.72106481]
[0.12114198 0.0788966  0.79996142]
[0.08686986 0.05658436 0.85654578]
[0.06229638 0.04057892 0.8971247 ]
[0.0446745 0.0291004 0.9262251]
[0.03203738 0.02086876 0.94709386]

```
\end{md}


To test the documentation gives the results that are written, run the following
at the command line:

\begin{cliin}
$ python -m doctest README.md
\end{cliin}


When testing the documentation, as for the testing of code, if there are no
errors there will be no output as shown in
Figure~\ref{fig:running_doctests_with_no_errors}.

\begin{figure}[htbp]
\centering

\includegraphics[width=0.750\linewidth]{./assets/running_doctests_with_no_errors/main.png}
\caption{Running the doctests with no errors.}
\label{fig:running_doctests_with_no_errors}
\end{figure}


Similarly to testing the code, if an error is included in the documentation an
error will be displayed when running the doc tests. This is shown in
Figure~\ref{fig:running_doctests_with_an_error} where the final output has been
changed to include an error: \texttt{-1} is written instead of \texttt{0.94709386}.

\begin{figure}[htbp]
\centering

\includegraphics[width=0.750\linewidth]{./assets/running_doctests_with_an_error/main.png}
\caption{Running the doctests with an error.}
\label{fig:running_doctests_with_an_error}
\end{figure}

Here is the fully modified tutorial of the documentation:

\begin{md}
# Absorption

Functionality to study the absorbing Markov chains.

## Tutorial

In this tutorial we will see how to use `absorption` to study an absorbing
Markov chain. For some background information on absorbing Markov chains we
recommend: <https://en.wikipedia.org/wiki/Absorbing_Markov_chain>.

Given a transition matrix $P$ defined by:

$$
p = \begin{pmatrix}
   1/2 & 1/4 & 1/4\\
   1/3 & 1/3 & 1/3\\
   0   & 0   & 1
    \end{pmatrix}
$$

We will start by seeing how the chain evolves over time by starting with an
initial vector $\pi=(1,0,0)$. In the next code snippet we will import the
necessary libraries and create both $P$ and $\pi$:

```python
>>> import numpy as np
>>> import absorption
>>> pi = np.array([1, 0, 0])
>>> P = np.array([[1 / 2, 1 / 4, 1 / 4], [1 / 3, 1 / 3, 1 / 3], [0, 0, 1]])

```

We now see how the vector $\pi$ changes over time:

```python
>>> for k in range(10):
...     print(absorption.get_long_run_state(pi, k, P))
[1. 0. 0.]
[0.5  0.25 0.25]
[0.33333333 0.20833333 0.45833333]
[0.23611111 0.15277778 0.61111111]
[0.16898148 0.1099537  0.72106481]
[0.12114198 0.0788966  0.79996142]
[0.08686986 0.05658436 0.85654578]
[0.06229638 0.04057892 0.8971247 ]
[0.0446745 0.0291004 0.9262251]
[0.03203738 0.02086876 0.94709386]

```

We see that, as expected over time the probability of being in the third state,
which is absorbing, increases.

We can also use `absorption` to get the average number of steps until
absorption from each state:

```python
>>> absorption.compute_t(P)
array([3.66666667, 3.33333333])

```

We see that the expected amounts of steps from the first state is slightly more
than from the second.
\end{md}

\subsection{Documenting the tests}
\label{\detokenize{building-tools/07-testing/tutorial/main:documenting-the-tests}}

Finally it is important to document how to run the tests. The \textbf{reference}
section is an appropriate place to put this. Add the following to the
\texttt{README.md} file:

\begin{md}
### Testing the software

To test the code:

```
$ python test_absorption.py
```

To test the documentation:

```
$ python -m doctest README.md
```
\end{md}

\section{How to}

\subsection{Write an \texttt{assert} statement}
\label{\detokenize{building-tools/07-testing/how/main:how-to-write-an-assert-statement}}

An \texttt{assert} statement is followed by 2 values: a boolean and an optional
string that gets returned if the boolean is \texttt{False}.


\begin{api}
assert boolean, string
\end{api}



For example, given a function that adds one to a variable:




\begin{pyin}
def add_one(a):
    """
    Returns a + 1
    """
    return a + 1
\end{pyin}

You can assert the expected behaviour:

\begin{pyin}
assert add_one(5) == 6, "The function gave the wrong answer."
\end{pyin}

Note that if you change the function to include an error for example here adding
2 and not 1, and run the same assert
you get an error as well as the specified string.

\begin{pyin}
def add_one(a):
    """
    Returns a + 1
    """
    return a + 2


assert add_one(5) == 6, "The function gave the wrong answer."
\end{pyin}





\begin{raw}
---------------------------------------------------------------------------
AssertionError                            Traceback (most recent call last)
Cell In[3], line 8
      2     """
      3     Returns a + 1
      4     """
      5     return a + 2
----> 8 assert add_one(5) == 6, "The function gave the wrong answer."

AssertionError: The function gave the wrong answer.
\end{raw}






\subsection{Write \texttt{assert} statements for code that acts randomly.}

When making an assertion about code that behaves in a random manner, use
seeding as described in Section~\ref{sec:reproduce_random_events}.


For example:




\begin{pyin}
import random


def roll_a_dice():
    """
    Pick a random integer between 1 and 6 (inclusive)
    """
    return random.choice(range(1, 7))
\end{pyin}





To test this, include a number of seeded assertions:




\begin{pyin}
random.seed(0)
assert roll_a_dice() == 4, "The 0 seed did not give the expected result"
random.seed(1)
assert roll_a_dice() == 2, "The 1 seed did not give the expected result"
random.seed(2)
assert roll_a_dice() == 1, "The 2 seed did not give the expected result"
random.seed(3)
assert roll_a_dice() == 2, "The 3 seed did not give the expected result"
\end{pyin}





You can also check behaviour over a number of repetitions:





\begin{pyin}
random.seed(0)
samples = [roll_a_dice() for repetition in range(1000)]
all_values = {1, 2, 3, 4, 5, 6}
assert set(samples) == all_values, "Not all values have been obtained over 1000 repetitions"
\end{pyin}






You can also confirm that the count of a given value is as expected:


\begin{pyin}
assert [samples.count(k) for k in range(1, 7)] == [193, 150, 166, 170, 152, 169], "The count of the values is not giving the expected count"
\end{pyin}


\begin{note}
The last assertion used the \texttt{count} method on a list that counts a given number
of items in a list.
\end{note}



\subsection{Write a test file.}

To write tests assertion statements should be put in to a file separate to the code
in functions.


For example, if the \texttt{dice.py} file contained:

\begin{python}
import random


def roll_a_dice():
    """
    Pick a random integer between 1 and 6 (inclusive)
    """
    return random.choice(range(1, 7))
\end{python}


Then a separate \texttt{test\_dice.py} file with the following would be written:

\begin{python}
import dice


def test_roll_a_dice_with_specific_values():
    """
    Check the roll a dice function gives specific numbers for a number of seeds.
    """
    random.seed(0)
    assert dice.roll_a_dice() == 4, "The 0 seed did not give the expected result"
    random.seed(1)
    assert dice.roll_a_dice() == 2, "The 1 seed did not give the expected result"
    random.seed(2)
    assert dice.roll_a_dice() == 1, "The 2 seed did not give the expected result"
    random.seed(3)
    assert dice.roll_a_dice() == 2, "The 3 seed did not give the expected result"


def test_roll_a_dice_for_a_large_sample():
    """
    Collect a sample of 1000 rolls and confirm that we have expected results.
    """
    random.seed(0)
    samples = [dice.roll_a_dice() for repetition in range(1000)]
    all_values = {1, 2, 3, 4, 5, 6}
    assert set(samples) == all_values, "Not all values have been obtained over 1000 repetitions"
    expected_counts = [193, 150, 166, 170, 152, 169]
    assert [samples.count(k) for k in range(1, 7)] == expected_counts, "The count of the values is not giving the expected count"

test_roll_a_dice_with_specific_values()
test_roll_a_dice_for_a_large_sample()
\end{python}


To run the tests you would then type the following at the command line:

\begin{cliin}
$ python test_dice.py
\end{cliin}


\subsection{Format doc tests.}

When writing code in documentation if you write it using a specific format then
python can be used to confirm that the code and its output is correct.


\begin{md}
>>> <python_code>
<expected_output
>>> <python_code_over_multiples_lines>
... <python_code_over_multiple_lines>
<expected_output>
\end{md}

\begin{itemize}
\item 

\texttt{>>>} is marks the start of a python statement.

\item 

\texttt{...} is used if the statement is multiple lines.

\item 

The expected output is written after the python statements.

\end{itemize}


For example if you were documenting the following code written in a \texttt{dice.py}
file:

\begin{python}
import random


def roll_a_dice():
    """
    Pick a random integer between 1 and 6 (inclusive)
    """
    return random.choice(range(1, 7))
\end{python}


You would write:

\begin{md}
>>> import dice
>>> random.seed(0)
>>> dice.roll_a_dice()
4

\end{md}


\subsection{Run doctests}
\label{\detokenize{building-tools/07-testing/how/main:how-to-run-doctests}}

Given a file with doc tests, to run them type the following at the
command line:


\begin{api}
$ python -m doctest <file>
\end{api}



For example:

\begin{cliin}
$ python -m doctest README.md
\end{cliin}




\section{Exercises}
\label{\detokenize{building-tools/07-testing/exercises/main:exercises}}\label{\detokenize{building-tools/07-testing/exercises/main::doc}}

Write tests for the \texttt{statistics.py} file written in the exercises of
Chapters~\ref{chp:modularisation} and~\ref{chp:documentation}.


Run the tests.



\section{Further information}
\label{\detokenize{building-tools/07-testing/why/main:further-information}}\label{\detokenize{building-tools/07-testing/why/main::doc}}

\subsection{Why are tests written as functions?}

In Section~\ref{sec:testing_tutorial} you wrote all the tests using
\texttt{assert} statements inside of functions. \textbf{Technically} this is not
necessary as you could write a single script with all the assert statements one
after the other.

This is not recommended: by using functions you directly have a place to document
the test (in the docstring) and the tests themselves are modularised.
Furthermore, this is actually how to write the tests when using a more
appropriate way of running tests as described in
Section~\ref{sec:is_there_a_more_efficient_way_to_run_tests}.


\subsection{Is there a more efficient way to run tests?}
\label{sec:is_there_a_more_efficient_way_to_run_tests}

Writing tests as a script and directly running them has one immediate problem:
once the first \texttt{assert} statement fails the rest of them are not run.


There is a Python library for running tests called
\texttt{pytest}~\cite{oliveira2018pytest}.

\subsection{What should be tested?}
\label{\detokenize{building-tools/07-testing/why/main:what-should-be-tested}}

The short answer to this is that all software should be tested and that software
is compromised of documentation and code.


Note that it is often not sufficient to test a function in a single way. For
example, consider a function that does two different things depending on the
parity of some input:

\begin{python}
def feedback_on_guess(guess, chosen_number):
    """
    Returns whether or not a guess is:

    - Larger than  a chosen_number
    - Smaller than a chosen_number
    - Equal to a chosen number
    """
    if guess < chosen_number:
        return "Guess is lower than chosen number"
    if guess > chosen_number:
        return "Guess is larger than chosen number"
    return "Guess is equal to chosen number"
\end{python}

In this case you would need to write at least 3 tests that check the 3
behaviours.
In practice there might be some functionality that is not tested but this should
be made clear and explicit and documented as to why should be written.


\subsection{Why do you need doc tests?}
\label{\detokenize{building-tools/07-testing/why/main:why-do-we-need-doc-tests}}

The purpose of doctests is to ensure that the code written in documentation is
correct.

It is important to not use doctests to test the functionality of the code: this
risks making the documentation unclear.


\subsection{What is test driven development?}
\label{\detokenize{building-tools/07-testing/why/main:what-is-test-driven-development}}

Test driven development is the development process of writing the test before
you write the code. Whilst this might seem counter-intuitive it is in fact a
strong approach to writing robust code efficiently.


In practice the process is as follows:
\begin{enumerate}

\item 

Write a test for some new functionality.

\item 

Run the tests to confirm that it fails (as the functionality is not yet
written).

\item 

Write the functionality.

\item 

Run the test.

\item 

Modify the test and/or the functionality

\end{enumerate}


Steps 4 and 5 might be repeated many times.


A good overview of test driven development is given in~\cite{percival2014test}.


\subsection{How are modularisation, documentation and testing related.}
\label{\detokenize{building-tools/07-testing/why/main:how-are-modularisation-documentation-and-testing-related}}

In Chapters~\ref{chp:modularisation},~\ref{chp:documentation}
and~\ref{chp:testing}
three concepts that move from writing code that works to writing software
that is reliable have been discussed::
\begin{itemize}
\item 

Modularisation.

\item 

Documentation.

\item 

Testing.

\end{itemize}


In reality \textbf{all three} of these concepts are closely related and fundamental to
good software. Figure\ref{fig:best_practice_triangle} shows this.

\begin{figure}[!htbp]
\centering

\includegraphics[width=0.750\linewidth]{./assets/best_practice_triangle/main.pdf}
\caption{The relationship between modularisation, documentation and testing}
\label{fig:best_practice_triangle}
\end{figure}
